\documentclass[a4paper,10pt]{article}
\usepackage[francais]{babel}
\usepackage{url}

\selectlanguage{francais}
\input{letter-ensiie.sty}

\selectlanguage{francais}
\enCouleur

\begin{document}


\name{St\'ephane \textsc{Genaud}}
\fonction{Directeur de l'antenne de Strasbourg de l'ENSIIE\\Professeur des Universités, ENSIIE}
\telephoneillkirch{+33 (0)3688 54542}
\email{genaud@ensiie.fr}

\begin{letter}[Avis concernant la candidature de Yufei WANG à la 2e année du
  diplôme d'Ingénieur Réseaux et Télécoms porté par Télécom Physique Strasbourg]{à}%
{Responsable du concours}

Chers Collègues,

Yufei WANG a intégré l'ENSIIE à la  rentrée de l'année scolaire 2015/2016 par la
voie  de l'admission  sur titre  dans  le cadre  des accords  de recrutement  de
l'Institut Mines-Télécom dans des universités partenaires Chinoises.

Yufei  a   effectué  sa  première  année   de  tronc  commun  à   Strasbourg  en
2015/2016. Après un premier semestre difficile, qui nécessite pour les étudiants
chinois un  effort d'adaptation culturel  et linguistique tout à  fait remarque,
Yufei a su prendre ses marques  dans notre environnement. Cette difficulté s'est
traduite  par un  semestre  d'automne non  validé.  Elle  a  validé au  semestre
d'automne 2015 les  UE Probabilités (4 ECTS), Graphes et  Optimisation (4 ECTS),
Formation Humaine (3 ECTS) et Economie-Gestion (3 ECTS). Elle a ensuite validé
son deuxième semestre en 2015/2016. Elle a redoublé son premier semestre cette
année scolaire 2016/2017 en validant les UEs défaillantes de l'année précédente,
qui concernaient essentiellement l'informatique. Les notes obtenues cette année
dans ces UEs montrent qu'elle a acquis un niveau satisfaisant (par exemple 14,5
en programmation impérative, 15 en bases de données, 10,5 en système).

Cette progression montre  que Yufei s'est bien adaptée à  un environnement comme
celui  de nos  écoles  d'ingénieur  françaises. Yufei  fait  preuve de  beaucoup
d'abnégation  dans le  travail  et  elle a  grandement  amélioré  sa capacité  à
travailler avec efficacité  vis à vis du mode d'évaluation  de notre système. Je
crois  qu'elle  pourra   également compter   sur  les  relations
personnelles  qu'elle a  nouées avec des étudiants de Télécom Physique Strasbourg
qu'elle cotoyait dans ce tronc commun, et c'est un élément social important pour
le bon déroulement de la scolarité. 

Je suis convaincu qu'elle sera à même de mener la scolarité proposée par le diplôme
pour lequel elle fait acte de candidature, et je recommande donc sa candidature.



\end{letter}
\begin{flushright}
\includegraphics[width=.2\textwidth]{signgenaud.jpg}
\end{flushright}
\end{document}
